\chapter{Construire un mini terminal, de A à Z}

Un terminal semble simple : on tape une commande, on lit la réponse. C'est
trompeur. Il doit lancer un programme, lui parler \emph{pendant} qu'il tourne,
comprendre un langage de contrôle vieux de cinquante ans, et ne jamais figer
l'interface --- sur trois systèmes qui ne s'y prennent pas du tout de la même
façon.

Dix étapes, en trois niveaux qui se justifient l'un l'autre. À la fin de
l'étape 3 vous avez quelque chose d'utile ; à la fin de l'étape 10, un vrai
terminal.

\begin{center}
\begin{tabular}{@{}llr@{}}
\toprule
\textbf{Fichier} & \textbf{Rôle} & \textbf{Lignes} \\
\midrule
\texttt{NkProcess.h} & lancer une commande, lire sa sortie & 171 \\
\texttt{NkPty.\{h,cpp\}} & tenir un shell \emph{interactif} vivant & 496 \\
\texttt{NkTerm.h} & interpréter le flux, en faire une grille & 631 \\
\bottomrule
\end{tabular}
\end{center}

\begin{nkretenir}
\textbf{Les étapes 1 et 2 du chapitre 21 ne changent pas.} Ajoutez
\texttt{NKThreading} aux dépendances --- il y est déjà --- et créez
\texttt{src/Terminal/}.
\end{nkretenir}

\section*{Niveau 1 --- lancer une commande sans geler}
\addcontentsline{toc}{section}{Niveau 1 --- lancer une commande sans geler}

\section{Étape 1 --- le fil d'arrière-plan et la file}

\begin{nkcode}{Applications/NKCode/src/NKCode/Project/NkProcess.h}
//   Start(cmd) lance la commande sur un THREAD d'arriere-plan (stdout+stderr
//   fusionnes) ; le thread pousse chaque ligne dans une file THREAD-SAFE ;
//   l'UI appelle Drain() chaque frame pour recuperer les nouvelles lignes sans
//   geler.
\end{nkcode}

\begin{nkcode}{La forme, en trois methodes}
class NkProcess {
    public:
        bool Start(const NkString &cmd);
        void Drain(NkVector<NkString> &out);   // vide la file, rend la main
        bool Running() const;
    private:
        void FilLecteur();                     // tourne dans un AUTRE fil
        NkMutex mMutex;
        NkVector<NkString> mLignes;
};
\end{nkcode}

\begin{nkretenir}
\textbf{C'est le patron « producteur lent, interface fluide », et il vaut bien
au-delà des terminaux.}

Un fil d'arrière-plan lit le programme et pousse ses lignes dans une file
protégée par un mutex. L'interface, à chaque image, appelle \texttt{Drain} : elle
prend ce qui est arrivé, ou rien, et rend la main immédiatement.

Elle n'\emph{attend} jamais. Une compilation de trois minutes défile ligne à
ligne pendant que la fenêtre reste vivante.
\end{nkretenir}

\begin{nkcode}{Applications/NKCode/src/NKCode/Project/NkProcess.h}
void Drain(NkVector<NkString> &out) {
    threading::NkScopedLock<NkMutex> lk(mMutex);
    for (/* ... */) { out.PushBack(mLines[i]); }
    // ... la file est videe
}
\end{nkcode}

\begin{nkpiege}
\textbf{Le verrou est pris pendant la copie, pas pendant la lecture du
programme.} C'est la règle du chapitre NKThreading appliquée : \emph{verrouiller
le temps de toucher la donnée, pas le temps de travailler}.

Si le fil gardait le verrou pendant qu'il attend la ligne suivante --- ce qui peut
durer des secondes --- l'interface bloquerait à chaque image en essayant de le
prendre. Vous auriez payé un fil pour obtenir un programme séquentiel qui
saccade.
\end{nkpiege}

\begin{nkretenir}[Une exception à la règle du dépôt, et elle est écrite]
\texttt{NkProcess.h} inclut \texttt{<cstdio>}, alors que le dépôt interdit la
bibliothèque standard. Son en-tête le dit et le justifie :

\begin{quote}
« WRAPPER MAISON DÉSIGNÉ pour les PIPES process : les appels
\texttt{\_popen}/\texttt{fgets}/\texttt{fgetc} de bas niveau sont volontairement
conservés dans cette famille de wrappers (pas d'équivalent \texttt{NkFile} pour
un flux de processus). »
\end{quote}

Comparez au lecteur vidéo du chapitre 23, qui incluait les mêmes en-têtes
\emph{sans rien dire}. Le code est identique ; ce qui diffère est décisif.

\textbf{Une exception écrite est une décision --- on peut la discuter, la lever,
la dater. Une exception muette est une dette que le prochain lecteur prendra
pour un exemple à suivre.}
\end{nkretenir}

\section{Étape 2 --- afficher, et savoir s'arrêter}

\begin{nkcode}{Dans OnTick}
void OnTick(float32) override {
    NkVector<NkString> arrivees;
    mProcessus.Drain(arrivees);
    for (uint64 i = 0; i < arrivees.Size(); ++i) { mLignes.PushBack(arrivees[i]); }
    if (mLignes.Size() > kMaxLignes) { /* jeter les plus anciennes */ }
}
\end{nkcode}

\begin{nkpiege}
\textbf{Bornez l'historique dès la première version.} Une commande bavarde
--- une compilation complète, un \texttt{find /} --- produit des centaines de
milliers de lignes. Sans plafond, la mémoire monte jusqu'à ce que le système
tue le programme.

Le symptôme est « le terminal a planté au bout de dix minutes », et la cause est
à l'endroit qu'on n'a pas écrit.
\end{nkpiege}

\begin{nkretenir}
\textbf{\texttt{NkProcess} lance une commande, la lit, et meurt.} C'est déjà
utile --- un bouton « Construire » qui affiche la sortie en direct --- et c'est
tout ce dont une partie des applications ont besoin.

Ne passez au niveau 2 que si vous voulez un \emph{vrai} terminal.
\end{nkretenir}

\section*{Niveau 2 --- un shell qui reste vivant}
\addcontentsline{toc}{section}{Niveau 2 --- un shell qui reste vivant}

\section{Étape 3 --- pourquoi un tube ne suffit pas}

Un \texttt{cd} lancé par \texttt{NkProcess} ne change rien : le processus meurt,
et le suivant repart du même dossier. Les variables aussi s'oublient. Il faut
garder \emph{le même} shell.

\begin{nkcode}{Applications/NKCode/src/NKCode/Project/NkPty.h}
// Contrairement a NkProcess (un _popen jetable par commande), NkPty lance UN
// shell (powershell / wsl -d <distro> / bash / cmd) attache a un pseudo-console
// Windows (ConPTY). On ECRIT les frappes clavier dans son entree et on LIT en
// continu sa sortie (avec sequences VT : couleurs, deplacement curseur, clear).
// => le shell affiche lui-meme son invite ; `clear` efface ; pas de boite de
//    saisie separee.
\end{nkcode}

\begin{nkretenir}
\textbf{« Le shell affiche lui-même son invite » est le renversement complet.}

Le débutant écrit une zone de saisie, un bouton \emph{Exécuter}, et une zone de
sortie. C'est une \emph{imitation} de terminal : l'invite est fausse, les
couleurs sont perdues, \texttt{vim} ne marche pas, \texttt{clear} n'efface rien.

Un vrai terminal ne fabrique rien. Il transmet les frappes au shell et affiche ce
que le shell renvoie --- \emph{y compris son invite}. Il n'est pas l'auteur de ce
qu'on voit ; il en est le verre.
\end{nkretenir}

\section{Étape 4 --- le pseudo-terminal}

Un programme se comporte différemment selon qu'il écrit dans un fichier ou vers
un humain : vers un fichier, il retire les couleurs et l'interactivité. Un tube
ordinaire ressemble à un fichier --- d'où des sorties ternes et \texttt{vim} qui
refuse de démarrer.

Un \textbf{pseudo-terminal} est un tube qui se \emph{présente} comme un
terminal : il a une taille en lignes et colonnes, il accepte d'être
redimensionné, et le programme d'en face le croit.

\begin{nkcode}{Applications/NKCode/src/NKCode/Project/NkPty.h}
bool Start(const NkString &cmdline, int16 cols, int16 rows,
           const NkString &cwd = NkString());
void Write(const char *data, usize len);   // les frappes clavier
void Write(const char *s);
void Resize(int16 cols, int16 rows);       // quand le panneau change de taille
\end{nkcode}

\begin{nkretenir}
\textbf{Quatre fonctions. C'est toute l'interface} --- et c'est ce qui rendra
l'étape 5 possible.
\end{nkretenir}

\begin{nkpiege}
\textbf{\texttt{Resize} n'est pas cosmétique.} Le programme d'en face décide sa
mise en page d'après la taille qu'on lui a déclarée. Oubliez de l'informer, et
\texttt{htop} dessinera pour 80 colonnes dans une fenêtre qui en fait 200 --- ou
débordera dans l'autre sens.

Un terminal qui ne réagit pas au redimensionnement n'est pas « moins abouti » :
il est faux dès la première fenêtre de taille inhabituelle.
\end{nkpiege}

\section{Étape 5 --- isoler le code de plateforme}

\begin{nkcode}{Applications/NKCode/src/NKCode/Project/NkPty.cpp (structure)}
#if defined(_WIN32)
    // ConPTY : CreatePseudoConsole, PROC_THREAD_ATTRIBUTE_PSEUDOCONSOLE
#else
  #if defined(__APPLE__)
    // <util.h>
  #else
    // <pty.h>   -- forkpty
  #endif
#endif
\end{nkcode}

\begin{nkretenir}
\textbf{Trois systèmes, trois mécanismes, un seul en-tête.} \texttt{NkPty.h}
n'expose que \texttt{Start}, \texttt{Write}, \texttt{Read}, \texttt{Resize} :
rien de ce qui suit ne connaît ConPTY ni \texttt{forkpty}.

L'en-tête dit d'ailleurs \emph{pourquoi} le \texttt{.cpp} existe : « isolée dans
\texttt{NkPty.cpp} pour ne pas polluer les en-têtes GUI avec les macros de
\texttt{windows.h} ». Ce n'est pas une préférence de style ---
\texttt{windows.h} définit des macros comme \texttt{min}, \texttt{max} et
\texttt{None} qui cassent le code qui les inclut par accident.

\textbf{Là où le portage est le plus sale, l'interface doit être la plus
étroite.}
\end{nkretenir}

\begin{nkpiege}
\textbf{\texttt{None} n'est pas une macro théorique.} X11 la définit à \texttt{0},
et NKGui l'emploie comme énumérateur --- ce qui a empêché NKGui de compiler sous
Linux pendant des mois. Le correctif a dû être posé \emph{à la porte d'entrée du
module}, parce qu'un \texttt{undef} par fichier ne protège pas ceux qui
l'incluent.
\end{nkpiege}

\section*{Niveau 3 --- comprendre ce que le shell renvoie}
\addcontentsline{toc}{section}{Niveau 3 --- comprendre ce que le shell renvoie}

\section{Étape 6 --- la grille de cellules}

Le shell ne renvoie pas du texte : il renvoie du texte \emph{mêlé d'ordres}.
\texttt{ESC[31m} veut dire « à partir d'ici, écris en rouge », \texttt{ESC[2J}
« efface l'écran », \texttt{ESC[10;5H} « place le curseur ligne 10, colonne 5 ».

\begin{nkcode}{Applications/NKCode/src/NKCode/Project/NkTerm.h}
struct NkTermCell {
        uint32 cp = 0x20;                  // codepoint Unicode (espace)
        NkColor fg = {204, 204, 204, 255};
        NkColor bg = {0, 0, 0, 0};         // a==0 => fond par defaut
};

class NkTerm {
    public:
        using Line = NkVector<NkTermCell>;
        NkTerm() { Resize(80, 24); }
        void Resize(int16 cols, int16 rows);
        // ... Feed(octets), TotalLines(), LineAt(i)
};
\end{nkcode}

\begin{nkretenir}
\textbf{L'émulateur transforme un flux en état.} Il consomme des octets et
maintient une grille : à chaque case, un caractère et deux couleurs. Le dessin,
ensuite, est trivial --- une boucle sur des cellules, en police à chasse fixe.

\textbf{C'est la séparation qui rend l'affaire faisable} : d'un côté un automate
qui ne dessine rien et \emph{se teste sans fenêtre} ; de l'autre un dessin qui ne
comprend rien et ne peut donc pas se tromper.

Vous reconnaissez le découpage règles/écran du chapitre 21, appliqué à un
terminal.
\end{nkretenir}

\begin{nkretenir}
Notez \texttt{bg = \{0,0,0,0\}} avec le commentaire « \texttt{a==0}
$\Rightarrow$ fond par défaut ». Une cellule ne dit pas « fond noir » : elle dit
« je n'ai pas d'avis ».

C'est ce qui permet au terminal de suivre le thème de l'éditeur. \textbf{Une
valeur \emph{absente} et une valeur \emph{noire} sont deux choses différentes},
et les confondre est l'erreur qui rend une interface non thémable.
\end{nkretenir}

\section{Étape 7 --- l'automate, et son état persistant}

\begin{nkcode}{Les quatre etats, et ils suffisent pour commencer}
enum class Etat : uint8 {
    TEXTE,      // on recopie les caracteres dans la grille
    ESC,        // on vient de lire 0x1B
    CSI,        // on a lu ESC[ ; on accumule les parametres
    OSC         // on a lu ESC] ; on consomme jusqu'au terminateur
};
\end{nkcode}

\begin{nkpiege}
\textbf{Une séquence d'échappement peut arriver coupée en deux.} La lecture rend
ce qui est disponible : \texttt{ESC[3} dans un paquet, \texttt{1m} dans le
suivant.

Un analyseur qui traite chaque paquet indépendamment affichera \texttt{[3} puis
\texttt{1m} en clair, aléatoirement, une fois de temps en temps --- le défaut qui
ne se reproduit jamais chez le développeur et toujours chez l'utilisateur.

\textbf{L'automate doit donc garder son état d'un appel à l'autre.} C'est le même
problème que le commentaire de bloc du chapitre 24 : \emph{un analyseur travaille
sur un flux, pas sur des morceaux}.
\end{nkpiege}

\begin{nkretenir}
\textbf{Ce qu'il faut interpréter en premier, et dans cet ordre} :
\texttt{SGR} (les couleurs), \texttt{ED} (effacer l'écran), \texttt{CUP} (placer
le curseur), puis les déplacements relatifs.

\textbf{Tout le reste se CONSOMME et s'ignore} --- mais se consomme. Une séquence
inconnue qu'on laisse passer s'affiche en clair et salit la sortie ; une séquence
inconnue qu'on avale proprement ne se voit pas.
\end{nkretenir}

\section{Étape 8 --- dessiner la grille}

\begin{nkcode}{La boucle, et elle est courte}
for (int32 l = 0; l < visibles; ++l) {
    const NkTerm::Line &ligne = mTerm.LineAt(premiere + l);
    for (int32 c = 0; c < mTerm.Cols(); ++c) {
        const NkTermCell &cel = ligne[c];
        if (cel.bg.a != 0) { dl.AddRectFilled(CaseRect(l, c), cel.bg); }
        if (cel.cp != 0x20) { dl.AddText(..., cel.cp, cel.fg); }
    }
}
\end{nkcode}

\begin{nkretenir}
\textbf{Une police à chasse fixe est obligatoire}, et ce n'est pas esthétique :
la grille suppose que toutes les cellules ont la même largeur. Avec une police
proportionnelle, les colonnes ne s'alignent plus et tout dessin ASCII se
disloque.

\textbf{Et l'on saute l'espace et le fond transparent} : sur un écran de terminal
typique, les trois quarts des cellules sont vides. Les dessiner quand même
multiplie le travail par quatre pour un résultat identique.
\end{nkretenir}

\section{Étape 9 --- le clavier, et le sens du flux}

\begin{nkcode}{Ce que fait la saisie -- et ce qu'elle ne fait PAS}
void OnKey(const NkKeyEvent &k) {
    // On ECRIT dans le shell. On ne touche JAMAIS a la grille.
    if (k.text) { mPty.Write(k.text); }
    else if (k.code == NK_ENTER)  { mPty.Write("\r"); }
    else if (k.code == NK_LEFT)   { mPty.Write("\x1b[D"); }
    // ...
}
\end{nkcode}

\begin{nkretenir}
\textbf{La saisie n'écrit jamais dans la grille.} Elle envoie au shell, le shell
renvoie l'écho, et \emph{c'est cet écho} qui apparaît.

Écrire soi-même le caractère tapé donnerait un double affichage --- le vôtre puis
celui du shell --- et casserait les mots de passe, que le shell n'écho justement
pas.

\textbf{Le flux va dans un seul sens : clavier $\rightarrow$ shell
$\rightarrow$ grille $\rightarrow$ écran.} Toute boucle courte est un défaut.
\end{nkretenir}

\begin{nkretenir}
Les touches de direction s'envoient comme des \emph{séquences} : \texttt{ESC[D}
pour la gauche. Le même langage sert dans les deux sens --- c'est pourquoi votre
automate de l'étape 7 vous aura appris à les écrire.
\end{nkretenir}

\section{Étape 10 --- redimensionner, et l'historique}

\begin{nkcode}{Dans OnLayout}
const int16 cols = (int16)(zone.w / mLargeurCar);
const int16 rows = (int16)(zone.h / mHauteurLigne);
if (cols != mTerm.Cols() || rows != mTerm.Rows()) {
    mTerm.Resize(cols, rows);   // la grille
    mPty.Resize(cols, rows);    // ET le shell -- les DEUX
}
\end{nkcode}

\begin{nkpiege}
\textbf{Les deux, toujours.} Redimensionner la grille sans prévenir le shell
donne un affichage juste sur des données calculées pour l'ancienne taille : les
lignes se replient au mauvais endroit et l'écran devient illisible après le
premier programme plein écran.
\end{nkpiege}

\begin{nkretenir}
\textbf{L'historique (\emph{scrollback}) se conserve au redimensionnement}, mais
l'écran se recadre. C'est ce que fait \texttt{NkTerm::Resize} du dépôt :
« conserve le scrollback ; recadre l'écran ».

Reflower tout l'historique à la nouvelle largeur est ce que font les terminaux
les plus aboutis, et c'est un chantier à soi seul --- ne le promettez pas dans
votre première version.
\end{nkretenir}

\begin{nkbilan}
Un terminal se construit en trois niveaux qui se justifient l'un l'autre :
\emph{lancer} sans geler (fil d'arrière-plan, file protégée, \texttt{Drain} par
image, historique borné) ; \emph{tenir un shell vivant} via un pseudo-terminal
--- c'est le shell qui affiche son invite, pas vous, et il faut lui déclarer sa
taille ; \emph{interpréter} le flux VT en une grille de cellules, avec un
automate qui garde son état d'un paquet à l'autre. Le code spécifique à chaque
système tient dans un seul \texttt{.cpp} derrière quatre fonctions. Et le flux va
dans un seul sens : ce que vous tapez part au shell, ce qui s'affiche revient de
lui.
\end{nkbilan}

\begin{nkquestions}
\begin{enumerate}
\item Pourquoi \texttt{Drain} plutôt qu'attendre la fin de la commande ?
\item Pourquoi faut-il borner l'historique dès la première version ?
\item Pourquoi un tube ordinaire ne suffit-il pas pour un shell interactif ?
\item Que se passe-t-il si l'on n'appelle jamais \texttt{Resize} --- sur la
      grille ? sur le shell ?
\item Pourquoi la saisie n'écrit-elle pas dans la grille ?
\item Pourquoi une cellule a-t-elle un fond d'alpha nul plutôt qu'un fond noir ?
\end{enumerate}
\end{nkquestions}

\begin{nkpratique}
Construisez votre mini terminal en suivant les dix étapes, dans l'ordre. Chaque
niveau doit se lancer avant que vous n'attaquiez le suivant.

\textbf{Niveau 1 (étapes 1--2).} Une zone de saisie, un bouton, une zone de
sortie. \emph{Vérifiez que l'interface reste fluide} pendant une commande de
trente secondes.

\textbf{Niveau 2 (étapes 3--5).} Un shell persistant. \texttt{cd} doit désormais
avoir un effet sur la commande suivante --- c'est votre critère de réussite, et
il est binaire.

\textbf{Niveau 3 (étapes 6--10).} Interprétez au minimum \texttt{SGR} et
\texttt{ED}. Le reste des séquences se consomme et s'ignore --- mais se consomme,
jamais ne s'affiche.
\end{nkpratique}

\begin{nkdemo}
Prouvez que votre analyseur survit à une séquence coupée.

Écrivez un banc \emph{sans fenêtre} --- vous pouvez, puisque
\texttt{NkTerm} ne dessine rien --- qui alimente votre émulateur avec
\texttt{"ESC[31mROUGE"} en un seul morceau, puis avec le même texte découpé
\textbf{à chaque octet}, un appel par caractère. Comparez la grille obtenue,
cellule par cellule.

Les deux doivent être identiques. Faites ensuite varier le point de coupure sur
\emph{toutes} les positions possibles et vérifiez les $n$ résultats.

\textbf{Pourquoi la variation exhaustive} : une coupure à un seul endroit peut
tomber juste par chance. C'est la coupure \emph{au milieu du nombre} ---
\texttt{ESC[3} puis \texttt{1m} --- qui casse les analyseurs naïfs, et rien ne
garantit que votre essai unique tombe dessus.
\end{nkdemo}

\begin{nklogique}
Votre terminal affiche parfois \texttt{[0m} en clair au milieu de la sortie ---
environ une fois sur cent lancements, jamais au même endroit.

\begin{enumerate}
\item Énumérez au moins trois causes possibles.
\item Pour chacune, dites quelle mesure la distinguerait des autres --- une
      mesure qui donne un résultat \emph{différent} selon la cause, sinon elle ne
      tranche rien.
\item Un collègue propose de filtrer \texttt{[0m} avant l'affichage. Expliquez
      pourquoi ce correctif est plus dangereux que le défaut, en une phrase.
\end{enumerate}
\end{nklogique}
